\documentclass{beamer}

% Theme choice
\usetheme{Madrid} % You can change to e.g., Warsaw, Berlin, CambridgeUS, etc.

% Encoding and font
\usepackage[utf8]{inputenc}
\usepackage[T1]{fontenc}

% Graphics and tables
\usepackage{graphicx}
\usepackage{booktabs}

% Code listings
\usepackage{listings}
\lstset{
basicstyle=\ttfamily\small,
keywordstyle=\color{blue},
commentstyle=\color{gray},
stringstyle=\color{red},
breaklines=true,
frame=single
}

% Math packages
\usepackage{amsmath}
\usepackage{amssymb}

% Colors
\usepackage{xcolor}

% TikZ and PGFPlots
\usepackage{tikz}
\usepackage{pgfplots}
\pgfplotsset{compat=1.18}
\usetikzlibrary{positioning}

% Hyperlinks
\usepackage{hyperref}

% Title information
\title{Chapter 3: Control Structures: Conditionals \& Loops}
\author{Your Name}
\institute{Your Institution}
\date{\today}

\begin{document}

\frame{\titlepage}

\begin{frame}[fragile]
    \frametitle{Introduction to Control Structures}

    Control structures are fundamental building blocks in programming that determine the flow of execution in a program. In Python, the primary control structures include conditionals and loops. 
\end{frame}

\begin{frame}[fragile]
    \frametitle{Importance of Control Structures}

    \begin{itemize}
        \item \textbf{Conditional Statements}: Enable decision-making within code, allowing a program to execute different blocks of code based on specific conditions.
            \begin{itemize}
                \item Example: Choosing an action based on user input (e.g., if a user is logged in or not).
            \end{itemize}
        
        \item \textbf{Loops}: Facilitate repetition of code execution, making it possible to automate tasks and manage collections of data efficiently.
            \begin{itemize}
                \item Example: Iterating through a list of items to compute totals or apply transformations.
            \end{itemize}
    \end{itemize}
\end{frame}

\begin{frame}[fragile]
    \frametitle{Code Snippet Example}

    \begin{lstlisting}[language=Python]
# Example of a Conditional Statement
user_age = 18

if user_age >= 18:
    print("You are eligible to vote.")
elif user_age > 0:
    print("You are not eligible to vote yet.")
else:
    print("Invalid age.")

# Example of a Loop
for number in range(1, 6):
    print(f"Current number: {number}")
    \end{lstlisting}
\end{frame}

\begin{frame}[fragile]
    \frametitle{Understanding Conditionals - Introduction}
    \begin{block}{What are Conditionals?}
        Conditionals are fundamental programming constructs that allow a program to make decisions based on specific conditions. They enable a program to execute certain blocks of code based on whether a condition is true or false.
    \end{block}
\end{frame}

\begin{frame}[fragile]
    \frametitle{Understanding Conditionals - Key Components}
    \begin{itemize}
        \item \textbf{if} - Checks a specified condition. If true, executes the block of code.
        \item \textbf{elif} - Allows checking multiple conditions. If any previous conditions are false, this checks the next condition.
        \item \textbf{else} - Executed when none of the previous conditions are true, providing a final alternative.
    \end{itemize}
\end{frame}

\begin{frame}[fragile]
    \frametitle{Understanding Conditionals - Syntax and Example}
    \begin{block}{Syntax}
        \begin{lstlisting}[language=Python]
if condition_1:
    # Code block for condition_1
elif condition_2:
    # Code block for condition_2
else:
    # Code block for when no conditions are met
        \end{lstlisting}
    \end{block}

    \begin{block}{Example Scenario}
        \begin{lstlisting}[language=Python]
score = 85

if score >= 90:
    grade = 'A'
elif score >= 80:
    grade = 'B'
elif score >= 70:
    grade = 'C'
else:
    grade = 'F'

print("The student's grade is:", grade)
        \end{lstlisting}
    \end{block}
\end{frame}

\begin{frame}[fragile]
    \frametitle{Understanding Conditionals - Key Points}
    \begin{itemize}
        \item \textbf{Indentation} - Crucial in Python as it defines the scope of code blocks controlled by conditionals.
        \item \textbf{True and False Evaluation} - Understanding how Python evaluates expressions (e.g., 0 is false, any non-zero number is true).
        \item \textbf{Multiple Conditions} - Use logical operators like \texttt{and}, \texttt{or}, and \texttt{not} to combine multiple conditions.
    \end{itemize}
\end{frame}

\begin{frame}[fragile]
    \frametitle{Understanding Conditionals - Conclusion}
    \begin{block}{Conclusion}
        Conditionals are vital for controlling the flow of a program based on dynamic conditions. Mastery of \texttt{if}, \texttt{elif}, and \texttt{else} statements is fundamental for implementing decision-making logic in Python programming.
    \end{block}
    \begin{block}{Next Steps}
        In the next slide, we will delve into practical examples illustrating these concepts in action.
    \end{block}
\end{frame}

\begin{frame}[fragile]
    \frametitle{Examples of Conditionals - Introduction}
    \begin{itemize}
        \item Conditional statements in Python enable code execution based on conditions.
        \item Key types of conditional statements:
        \begin{itemize}
            \item \textbf{if statement}: Executes if the condition is true.
            \item \textbf{elif statement}: Checks additional conditions.
            \item \textbf{else statement}: Executes if all previous conditions are false.
        \end{itemize}
    \end{itemize}
\end{frame}

\begin{frame}[fragile]
    \frametitle{Examples of Conditionals - General Syntax}
    \begin{block}{General Syntax}
        \begin{lstlisting}[language=Python]
if condition1:
    # Code to execute if condition1 is True
elif condition2:
    # Code to execute if condition2 is True
else:
    # Code to execute if all previous conditions are False
        \end{lstlisting}
    \end{block}
\end{frame}

\begin{frame}[fragile]
    \frametitle{Examples of Conditionals - Grading System}
    \begin{block}{Example 1: Grading System}
        \begin{lstlisting}[language=Python]
score = 85

if score >= 90:
    grade = 'A'
elif score >= 80:
    grade = 'B'
elif score >= 70:
    grade = 'C'
elif score >= 60:
    grade = 'D'
else:
    grade = 'F'

print("Your grade is:", grade)
        \end{lstlisting}
    \end{block}
    \begin{itemize}
        \item The \texttt{if} statement checks for scores 90 or above.
        \item \texttt{elif} checks other ranges for grades.
        \item The \texttt{else} assigns an F for scores below 60.
    \end{itemize}
\end{frame}

\begin{frame}[fragile]
    \frametitle{Examples of Conditionals - Age Verification}
    \begin{block}{Example 2: Age Verification}
        \begin{lstlisting}[language=Python]
age = 20

if age >= 18:
    print("You are eligible to vote.")
else:
    print("You are not eligible to vote.")
        \end{lstlisting}
    \end{block}
    \begin{itemize}
        \item Simple structure with \texttt{if} and \texttt{else}.
        \item Checks if age is 18 or older for voting eligibility.
    \end{itemize}
\end{frame}

\begin{frame}[fragile]
    \frametitle{Examples of Conditionals - Weather Comparison}
    \begin{block}{Example 3: Weather Comparison}
        \begin{lstlisting}[language=Python]
weather = "rainy"

if weather == "sunny":
    activity = "Go for a picnic."
elif weather == "cloudy":
    activity = "Visit a museum."
elif weather == "rainy":
    activity = "Stay indoors and read a book."
else:
    activity = "Check the weather and plan accordingly."

print("Suggested activity:", activity)
        \end{lstlisting}
    \end{block}
    \begin{itemize}
        \item Multiple conditions checked using \texttt{elif}.
        \item \texttt{else} handles unforeseen weather conditions.
    \end{itemize}
\end{frame}

\begin{frame}[fragile]
    \frametitle{Examples of Conditionals - Conclusion}
    \begin{itemize}
        \item Conditional statements enhance flexibility and control flow in programs.
        \item Useful for decision-making, validations, etc.
        \item Promote code readability and maintainability.
        \item Effective use allows management of complex logic in applications.
    \end{itemize}
\end{frame}

\begin{frame}
    \frametitle{Nested Conditionals}
    \begin{itemize}
        \item Nested conditionals are decision-making constructs used within another conditional.
        \item In Python, this is achieved by placing an \texttt{if} statement inside another \texttt{if} statement.
        \item Careless use can reduce code readability and increase complexity.
    \end{itemize}
\end{frame}

\begin{frame}[fragile]
    \frametitle{Understanding Nested Conditionals}
    \begin{block}{Definition}
        A nested conditional occurs when one conditional statement is placed inside another. The inner condition is only evaluated if the outer condition is true.
    \end{block}
    
    \begin{block}{Syntax Example}
    \begin{lstlisting}[language=Python]
if condition1:
    # Outer conditional
    if condition2:
        # Inner conditional
        print("Both conditions are true.")
    else:
        print("Only the outer condition is true.")
else:
    print("Outer condition is false.")
    \end{lstlisting}
    \end{block}
\end{frame}

\begin{frame}[fragile]
    \frametitle{Example Scenario: Grading System}
    Consider a scenario where a student’s grade is evaluated based on two conditions: their score and attendance rate.
    
    \begin{lstlisting}[language=Python]
score = 85
attendance = 90

if score >= 70:
    if attendance >= 75:
        print("Student is passing with good standing.")
    else:
        print("Student is passing but has low attendance.")
else:
    print("Student is failing.")
    \end{lstlisting}
    
    \begin{itemize}
        \item The outer condition checks if the score is passing (70 or higher).
        \item The inner condition checks if the attendance is adequate (75\% or higher).
    \end{itemize}
\end{frame}

\begin{frame}
    \frametitle{Implications for Readability and Complexity}
    \begin{block}{Key Points to Consider}
        \begin{itemize}
            \item \textbf{Readability:}
                \begin{itemize}
                    \item Excessive nesting can obscure code clarity and maintainability.
                    \item Strive for clarity; each nested level should serve a specific purpose.
                \end{itemize}
            \item \textbf{Complexity:}
                \begin{itemize}
                    \item Each extra layer of nesting increases cognitive load.
                    \item Simplify logic to improve maintainability.
                \end{itemize}
        \end{itemize}
    \end{block}
    
    \begin{block}{Best Practices}
        \begin{itemize}
            \item Limit nesting levels to 2-3 if possible.
            \item Utilize logical operators (\texttt{and}, \texttt{or}) to simplify conditions.
            \item Refactor complex structures into functions for clarity.
        \end{itemize}
    \end{block}
\end{frame}

\begin{frame}
    \frametitle{Conclusion}
    Nested conditionals are useful tools in programming that enable detailed decision-making. However, it is crucial to use them judiciously. 
    \begin{itemize}
        \item Strive for clarity and simplicity.
        \item Ensuring code remains readable and maintainable is essential for long-term software development.
    \end{itemize}
\end{frame}

\begin{frame}[fragile]
    \frametitle{Introduction to Loops - Overview}
    \begin{itemize}
        \item Loops are control structures for executing a block of code repeatedly.
        \item In Python, there are two primary types:
        \begin{itemize}
            \item \textbf{For loops}
            \item \textbf{While loops}
        \end{itemize}
        \item They help automate repetitive tasks, enhance efficiency, and reduce redundancy.
    \end{itemize}
\end{frame}

\begin{frame}[fragile]
    \frametitle{Types of Loops - For Loops}
    \begin{block}{For Loops}
        \begin{itemize}
            \item \textbf{Purpose}: To iterate over a sequence (e.g., list, tuple).
            \item \textbf{Syntax}:
            \begin{lstlisting}
for variable in sequence:
    # Code to execute
            \end{lstlisting}
            \item \textbf{Example}:
            \begin{lstlisting}
fruits = ["apple", "banana", "cherry"]
for fruit in fruits:
    print(fruit)
            \end{lstlisting}
            \item \textbf{Output}:
            \begin{lstlisting}
apple
banana
cherry
            \end{lstlisting}
        \end{itemize}
    \end{block}
\end{frame}

\begin{frame}[fragile]
    \frametitle{Types of Loops - While Loops}
    \begin{block}{While Loops}
        \begin{itemize}
            \item \textbf{Purpose}: Repeats a block of code while a condition is true.
            \item \textbf{Syntax}:
            \begin{lstlisting}
while condition:
    # Code to execute
            \end{lstlisting}
            \item \textbf{Example}:
            \begin{lstlisting}
count = 0
while count < 5:
    print(count)
    count += 1
            \end{lstlisting}
            \item \textbf{Output}:
            \begin{lstlisting}
0
1
2
3
4
            \end{lstlisting}
        \end{itemize}
    \end{block}
\end{frame}

\begin{frame}[fragile]
    \frametitle{Use Cases and Key Points}
    \begin{block}{Use Cases}
        \begin{itemize}
            \item \textbf{For Loops}:
            \begin{itemize}
                \item Iterating through collections and strings.
                \item Performing operations a set number of times.
            \end{itemize}
            \item \textbf{While Loops}:
            \begin{itemize}
                \item Repeating tasks until a condition changes.
                \item Implementing indefinite algorithms.
            \end{itemize}
        \end{itemize}
    \end{block}

    \begin{block}{Key Points}
        \begin{itemize}
            \item Use \textbf{for loops} when iterations are known.
            \item Use \textbf{while loops} for unspecified iterations.
            \item Beware of infinite loops, ensure conditions become false.
        \end{itemize}
    \end{block}
\end{frame}

\begin{frame}[fragile]
    \frametitle{For Loops - Introduction}
    \begin{itemize}
        \item A \textbf{for loop} in Python allows iterating over sequences (lists, tuples, strings).
        \item Simplifies tasks involving repeated operations based on collections.
    \end{itemize}
\end{frame}

\begin{frame}[fragile]
    \frametitle{For Loops - Syntax}
    \begin{block}{Syntax}
        The general syntax of a for loop is as follows:
        \begin{lstlisting}[language=Python]
for item in iterable:
    # Code block to execute
        \end{lstlisting}
    \end{block}
    \begin{itemize}
        \item \texttt{item}: A variable representing each element in the iterable.
        \item \texttt{iterable}: A collection (list, tuple, string) of items to iterate over.
    \end{itemize}
\end{frame}

\begin{frame}[fragile]
    \frametitle{For Loops - Examples}
    \begin{block}{Example 1: Iterating Over a List}
        \begin{lstlisting}[language=Python]
fruits = ['apple', 'banana', 'cherry']
for fruit in fruits:
    print(fruit)
        \end{lstlisting}
        \textbf{Output:}
        \begin{verbatim}
apple
banana
cherry
        \end{verbatim}
    \end{block}
    
    \begin{block}{Example 2: Iterating Over a Tuple}
        \begin{lstlisting}[language=Python]
colors = ('red', 'green', 'blue')
for color in colors:
    print(color)
        \end{lstlisting}
        \textbf{Output:}
        \begin{verbatim}
red
green
blue
        \end{verbatim}
    \end{block}
\end{frame}

\begin{frame}[fragile]
    \frametitle{For Loops - More Examples}
    \begin{block}{Example 3: Iterating Over a String}
        \begin{lstlisting}[language=Python]
message = "Hello"
for char in message:
    print(char)
        \end{lstlisting}
        \textbf{Output:}
        \begin{verbatim}
H
e
l
l
o
        \end{verbatim}
    \end{block}
\end{frame}

\begin{frame}[fragile]
    \frametitle{For Loops - Key Points}
    \begin{itemize}
        \item \textbf{Flexibility}: Works with any iterable data type.
        \item \textbf{Readability}: Intuitive and enhances code readability.
        \item \textbf{Efficiency}: Eliminates manual index management.
    \end{itemize}
\end{frame}

\begin{frame}[fragile]
    \frametitle{For Loops - Common Use Cases}
    \begin{enumerate}
        \item Processing \textbf{elements} in a collection.
        \item Modifying element \textbf{values} in lists/tuples.
        \item \textbf{Generating} new sequences based on existing data.
    \end{enumerate}
\end{frame}

\begin{frame}[fragile]
    \frametitle{For Loops - Conclusion}
    \begin{itemize}
        \item For loops streamline repetitive tasks in Python programming.
        \item Understanding their syntax and use cases leads to more efficient and readable code.
    \end{itemize}
\end{frame}

\begin{frame}[fragile]
    \frametitle{While Loops - Overview}
    \begin{itemize}
        \item A \textbf{while loop} repeats a block of code while a specified condition is true.
        \item Useful for dynamic repetition when the number of iterations is unknown.
        \item \textbf{Syntax:}
        \begin{lstlisting}[language=Python]
while condition:
    # block of code to be executed
        \end{lstlisting}
    \end{itemize}
\end{frame}

\begin{frame}[fragile]
    \frametitle{While Loops - Mechanics}
    \begin{enumerate}
        \item \textbf{Condition Evaluation}:
            \begin{itemize}
                \item Before each iteration, the condition is checked.
                \item If \textbf{true}, the code block executes.
                \item If \textbf{false}, the loop terminates.
            \end{itemize}
    \end{enumerate}
\end{frame}

\begin{frame}[fragile]
    \frametitle{While Loops - Example}
    \begin{block}{Example Code}
        \begin{lstlisting}[language=Python]
# Initialize a counter variable
counter = 0

# Start a while loop
while counter < 5:
    print("Counter is:", counter)
    counter += 1  # Increment the counter
        \end{lstlisting}
    \end{block}
    \begin{block}{Output}
        \begin{lstlisting}
Counter is: 0
Counter is: 1
Counter is: 2
Counter is: 3
Counter is: 4
        \end{lstlisting}
    \end{block}
    \begin{itemize}
        \item This loop continues until \texttt{counter} reaches 5.
    \end{itemize}
\end{frame}

\begin{frame}[fragile]
    \frametitle{Loop Control Statements}
    \begin{itemize}
        \item \textbf{Control Statements} manage flow in loops: \texttt{break} and \texttt{continue}.
    \end{itemize}
    
    \begin{block}{1. Break Statement}
        \begin{itemize}
            \item Exits the loop immediately, even if the condition is true.
            \item Useful for ending the loop under specific conditions.
        \end{itemize}
        \begin{lstlisting}[language=Python]
counter = 0

while True:  # Infinite loop
    if counter == 3:
        break  # Exit the loop when counter is 3
    print("Counter is:", counter)
    counter += 1
        \end{lstlisting}
    \end{block}

\end{frame}

\begin{frame}[fragile]
    \frametitle{Loop Control Statements - Continue}
    \begin{block}{2. Continue Statement}
        \begin{itemize}
            \item Skips the current iteration and moves to the next one.
            \item Useful to avoid executing remaining code in certain conditions.
        \end{itemize}
        \begin{lstlisting}[language=Python]
counter = 0

while counter < 5:
    counter += 1
    if counter == 3:
        continue  # Skip the print statement when counter is 3
    print("Counter is:", counter)
        \end{lstlisting}
    \end{block}
    \begin{block}{Output}
        \begin{lstlisting}
Counter is: 1
Counter is: 2
Counter is: 4
Counter is: 5
        \end{lstlisting}
    \end{block}
\end{frame}

\begin{frame}
    \frametitle{Key Points to Remember}
    \begin{itemize}
        \item Ensure the condition will eventually become false to avoid infinite loops.
        \item Use \texttt{break} and \texttt{continue} to control loop flow efficiently.
        \item Practical use cases include reading inputs until a sentinel value is received.
    \end{itemize}
\end{frame}

\begin{frame}[fragile]
    \frametitle{Loop Control Statements - Overview}
    \begin{itemize}
        \item Loop control statements are essential for managing the flow of execution within loops.
        \item Focus on three key statements:
            \begin{itemize}
                \item \texttt{break}
                \item \texttt{continue}
                \item \texttt{pass}
            \end{itemize}
        \item Understanding these statements improves loop efficiency and exit condition management.
    \end{itemize}
\end{frame}

\begin{frame}[fragile]
    \frametitle{Loop Control Statements - \texttt{break} Statement}
    \begin{block}{Explanation}
        The \texttt{break} statement is used to terminate a loop prematurely. Control transfers to the next statement following the loop.
    \end{block}
    
    \begin{block}{Use Cases}
        \begin{itemize}
            \item Exiting a loop when a certain condition is met (e.g., user input).
            \item Finding a target value.
        \end{itemize}
    \end{block}
    
    \begin{lstlisting}[language=Python]
for i in range(10):
    if i == 5:
        break  # Exit the loop when i equals 5
    print(i)
    \end{lstlisting}
    
    \begin{block}{Output}
        \begin{verbatim}
0
1
2
3
4
        \end{verbatim}
    \end{block}
\end{frame}

\begin{frame}[fragile]
    \frametitle{Loop Control Statements - \texttt{continue} and \texttt{pass}}
    \begin{block}{\texttt{continue} Statement}
        The \texttt{continue} statement skips the current iteration and proceeds to the next iteration of the loop.
    \end{block}
    
    \begin{block}{Use Cases}
        \begin{itemize}
            \item Skipping iterations based on specific conditions (e.g., ignoring negative numbers).
        \end{itemize}
    \end{block}
    
    \begin{lstlisting}[language=Python]
for i in range(10):
    if i % 2 == 0:
        continue  # Skip even numbers
    print(i)
    \end{lstlisting}
    
    \begin{block}{Output}
        \begin{verbatim}
1
3
5
7
9
        \end{verbatim}
    \end{block}
\end{frame}

\begin{frame}[fragile]
    \frametitle{Loop Control Statements - \texttt{pass} Statement}
    \begin{block}{Explanation}
        The \texttt{pass} statement is used as a placeholder; it does nothing but allows for syntactically correct code.
    \end{block}
    
    \begin{block}{Use Cases}
        \begin{itemize}
            \item Serves as a placeholder during development for future code implementation.
        \end{itemize}
    \end{block}
    
    \begin{lstlisting}[language=Python]
for i in range(5):
    if i < 3:
        pass  # Placeholder for future code
    else:
        print(i)
    \end{lstlisting}
    
    \begin{block}{Output}
        \begin{verbatim}
3
4
        \end{verbatim}
    \end{block}
\end{frame}

\begin{frame}[fragile]
    \frametitle{Key Points to Emphasize}
    \begin{itemize}
        \item \textbf{break}: Exits the loop entirely.
        \item \textbf{continue}: Skips the current iteration and continues with the next one.
        \item \textbf{pass}: Acts as a placeholder for future implementations.
    \end{itemize}
    
    Understanding and effectively using these control statements allows programmers to implement flexible and efficient loop logic, enhancing code design and clarity.
\end{frame}

\begin{frame}
    \frametitle{Combining Conditionals and Loops}
    \begin{itemize}
        \item Powerful programming technique for decision-making and repetition
        \item Used to solve complex real-world problems dynamically
    \end{itemize}
\end{frame}

\begin{frame}
    \frametitle{Concept Overview}
    \begin{itemize}
        \item \textbf{Conditionals}:
            \begin{itemize}
                \item Execute different sections of code based on conditions
                \item Common statements: \texttt{if}, \texttt{else if}, \texttt{else}
            \end{itemize}
        \item \textbf{Loops}:
            \begin{itemize}
                \item Repeat a block of code multiple times
                \item Common types: \texttt{for} loops and \texttt{while} loops
            \end{itemize}
    \end{itemize}
\end{frame}

\begin{frame}[fragile]
    \frametitle{Example 1: User Login System}
    \begin{lstlisting}[language=Python]
correct_password = "securePassword123"
attempts = 3

while attempts > 0:
    user_input = input("Enter your password: ")
    if user_input == correct_password:
        print("Login successful!")
        break
    else:
        attempts -= 1
        print(f"Incorrect password. You have {attempts} attempts left.")

if attempts == 0:
    print("Access denied.")
    \end{lstlisting}
    \begin{block}{Explanation}
        The loop continues until the user enters the correct password 
        or runs out of attempts. Conditionals check the password's 
        correctness, providing feedback accordingly.
    \end{block}
\end{frame}

\begin{frame}[fragile]
    \frametitle{Example 2: Grade Feedback System}
    \begin{lstlisting}[language=Python]
scores = [85, 92, 58, 75, 99]
feedback = []

for score in scores:
    if score >= 90:
        feedback.append("Excellent")
    elif score >= 75:
        feedback.append("Good")
    else:
        feedback.append("Needs Improvement")

print(feedback)
    \end{lstlisting}
    \begin{block}{Explanation}
        The loop iterates through a list of scores. Conditionals categorize 
        these scores into feedback labels, constructing a list of feedback.
    \end{block}
\end{frame}

\begin{frame}
    \frametitle{Key Points to Emphasize}
    \begin{enumerate}
        \item \textbf{Flexibility in Programming}:
            \begin{itemize}
                \item Allows handling various scenarios dynamically
            \end{itemize}
        \item \textbf{Efficiency}:
            \begin{itemize}
                \item Reduces redundancy in code through loops
            \end{itemize}
        \item \textbf{User Interaction}:
            \begin{itemize}
                \item Essential for interactive programs responding to user actions
            \end{itemize}
    \end{enumerate}
\end{frame}

\begin{frame}
    \frametitle{Best Practices}
    \begin{itemize}
        \item \textbf{Always Validate Input}:
            \begin{itemize}
                \item Ensure graceful handling of invalid inputs in loops
            \end{itemize}
        \item \textbf{Limit Loop Iterations}:
            \begin{itemize}
                \item Provide a clear exit condition to avoid infinite loops
            \end{itemize}
    \end{itemize}
\end{frame}

\begin{frame}[fragile]
    \frametitle{Best Practices for Control Structures}
    \begin{enumerate}
        \item \textbf{Keep it Simple and Clear}
        \begin{itemize}
            \item Aim for readable and understandable code.
            \item Prefer simple conditional checks over complex nested conditions.
            \item \textbf{Example}: 
            \begin{lstlisting}
if user_age >= 18:
    print("You are an adult.")
            \end{lstlisting}
        \end{itemize}

        \item \textbf{Use Meaningful Names}
        \begin{itemize}
            \item Use descriptive variable and function names to convey their purpose.
            \item \textbf{Example}: Instead of \(x\), use \(user\_age\).
        \end{itemize}

        \item \textbf{Avoid Deep Nesting}
        \begin{itemize}
            \item Keep control structures flat to enhance readability.
            \item Consider using \texttt{elif} instead of multiple nested \texttt{if} statements.
            \item \textbf{Example}:
            \begin{lstlisting}
if condition1:
    # do something
elif condition2:
    # do something else
            \end{lstlisting}
        \end{itemize}
    \end{enumerate}
\end{frame}

\begin{frame}[fragile]
    \frametitle{Best Practices for Control Structures (Contd.)}
    \begin{enumerate}[resume]
        \item \textbf{Minimize Side Effects}
        \begin{itemize}
            \item Ensure conditionals do not change the state of variables unintentionally.
            \item Use functions to encapsulate behavior with clear inputs and outputs.
        \end{itemize}

        \item \textbf{Document Your Code}
        \begin{itemize}
            \item Include comments to explain non-obvious logic or decisions.
            \item \textbf{Example}: 
            \begin{lstlisting}
# Check if user is a minor
if user_age < 18:
    print("You are a minor.")
            \end{lstlisting}
        \end{itemize}
    \end{enumerate}
\end{frame}

\begin{frame}[fragile]
    \frametitle{Common Mistakes to Avoid}
    \begin{enumerate}
        \item \textbf{Incorrect Condition Evaluation}
        \begin{itemize}
            \item Be cautious with equality and relational operators; mixing them up can create logic errors.
            \item \textbf{Example}: 
            \begin{lstlisting}
# Incorrect: 
if user_age = 18: # assignment instead of equality check
    print("You are exactly 18.")
            \end{lstlisting}
        \end{itemize}

        \item \textbf{Neglecting Edge Cases}
        \begin{itemize}
            \item Account for all possible input values, including unusual cases.
            \item \textbf{Example}: Not checking for invalid inputs (e.g., negative ages).
        \end{itemize}

        \item \textbf{Overusing Global Variables}
        \begin{itemize}
            \item Minimize global variable usage to prevent unpredictable behaviors.
            \item Instead, pass variables as parameters or use class attributes for state.
        \end{itemize}
    \end{enumerate}
\end{frame}

\begin{frame}[fragile]
    \frametitle{Common Mistakes to Avoid (Contd.)}
    \begin{enumerate}[resume]
        \item \textbf{Failing to Break Loops Appropriately}
        \begin{itemize}
            \item Ensure exit conditions are included in loops to avoid infinite loops.
            \item \textbf{Example}:
            \begin{lstlisting}
while True:
    if some_condition:
        break  # Ensure this condition is met to exit the loop
            \end{lstlisting}
        \end{itemize}

        \item \textbf{Lack of Consistent Indentation and Formatting}
        \begin{itemize}
            \item Consistent formatting aids readability and helps identify structures.
        \end{itemize}
    \end{enumerate}
\end{frame}

\begin{frame}[fragile]
    \frametitle{Key Takeaways}
    \begin{itemize}
        \item Strive for readability and maintainability in control structures.
        \item Use clear naming, minimize nesting, and avoid unintended side effects.
        \item Double-check conditions, handle edge cases, and ensure loop exits are accounted for.
        \item Comment your code thoughtfully to aid understanding and future revisions.
    \end{itemize}
\end{frame}


\end{document}